\documentclass{article}
\usepackage{amsmath,amssymb,amsthm}
\usepackage{enumitem}
\usepackage{hyperref}
\usepackage{xcolor}
\newtheorem{theorem}{Theorem}
\newtheorem{lemma}{Lemma}
\newcommand{\E}{\mathbb{E}}
\newcommand{\R}{\mathbb{R}}

\begin{document}

\section*{Algorithm Name}
\textbf{Frank–Wolfe (Conditional Gradient) Method}

\section*{Mathematical Formulation}
Let $f:\R^{d}\rightarrow\R$ be a convex and $L$‑smooth function and let 
\[
\mathcal C\subset\R^{d}
\]
be a non‑empty compact convex set (the feasible domain).  
Starting from an initial point $x_{0}\in\mathcal C$, the algorithm generates a sequence
$\{x_{k}\}_{k\ge 0}$ by the following two steps at iteration $k$:

\begin{enumerate}[label=\arabic*.]
\item \textbf{Linear minimization oracle (LMO)}  
\[
s_{k}\;=\;\arg\min_{s\in\mathcal C}\;\langle \nabla f(x_{k}),\,s\rangle .
\tag{1}
\]

\item \textbf{Convex combination update}  
\[
x_{k+1}\;=\;(1-\gamma_{k})\,x_{k}\;+\;\gamma_{k}\,s_{k},
\qquad \gamma_{k}\in[0,1].
\tag{2}
\]
\end{enumerate}
Typical step‑size choices are  

\begin{itemize}
\item \emph{Fixed}: $\gamma_{k}=\gamma\in(0,1]$,
\item \emph{Diminishing}: $\displaystyle \gamma_{k}=\frac{2}{k+2}$ (the classic Frank–Wolfe schedule),
\item \emph{Exact line‑search}: $\displaystyle 
\gamma_{k}=\arg\min_{\gamma\in[0,1]} f\big((1-\gamma)x_{k}+\gamma s_{k}\big)$.
\end{itemize}

\section*{Pseudocode}
\begin{verbatim}
Input:  convex smooth f, feasible set C, 
        initial point x0 ∈ C, 
        step‑size rule {γk}
Output: sequence {xk}

x ← x0
for k = 0,1,2,… do
    // 1) linear minimization oracle
    sk ← argmin_{s∈C} ⟨∇f(x), s⟩
    
    // 2) step‑size (choose one of the rules above)
    γ ← step_size(k)   // e.g. γ = 2/(k+2)
    
    // 3) convex combination update
    x ← (1‑γ) * x + γ * sk
    
    // optional termination test
    if stopping_criterion_met then break
end for
return x
\end{verbatim}

\section*{Dry Run Example}
Consider the one‑dimensional problem  

\[
f(w)=(w-3)^{2},\qquad \mathcal C=[0,5],\qquad w_{0}=0,
\]
and use the fixed step size $\gamma=0.1$.  
The gradient is $\nabla f(w)=2(w-3)$.  

\begin{center}
\begin{tabular}{c|c|c|c|c}
$k$ & $w_{k}$ & $\nabla f(w_{k})$ & $s_{k}=\arg\min_{s\in[0,5]}\langle\nabla f(w_{k}),s\rangle$ & $w_{k+1}=(1-\gamma)w_{k}+\gamma s_{k}$\\ \hline
0 & $0$ & $-6$ & $5$ (since $-6s$ is smallest at $s=5$) & $0.1\cdot5=0.5$\\
1 & $0.5$ & $-5$ & $5$ & $0.9\cdot0.5+0.1\cdot5=0.95$\\
2 & $0.95$ & $-4.1$ & $5$ & $0.9\cdot0.95+0.1\cdot5=1.355$\\
3 & $1.355$ & $-3.29$ & $5$ & $0.9\cdot1.355+0.1\cdot5=1.7195$
\end{tabular}
\end{center}
Corresponding objective values $f(w_{k})$ are $9$, $6.25$, $4.2025$, $2.852\ldots$, showing a monotone decrease.

\newpage

\section*{Assumptions}
\begin{enumerate}[label=\textbf{A\arabic*.},leftmargin=*]
\item \label{A1} \textbf{Convexity.} $f$ is convex on $\mathcal C$:  
\[
f(y)\ge f(x)+\langle\nabla f(x),y-x\rangle,\qquad\forall x,y\in\mathcal C .
\]

\item \label{A2} \textbf{$L$‑smoothness.} $f$ has $L$‑Lipschitz continuous gradient on $\mathcal C$:  
\[
\|\nabla f(x)-\nabla f(y)\|\le L\|x-y\|,\qquad\forall x,y\in\mathcal C .
\]

\item \label{A3} \textbf{Compactness of $\mathcal C$.}  
The set $\mathcal C$ is compact and convex with (Euclidean) diameter
\[
D\;:=\;\max_{x,y\in\mathcal C}\|x-y\|<\infty .
\]

\item \label{A4} \textbf{Existence of an optimal solution.}  
There exists $x^{*}\in\mathcal C$ such that $f(x^{*})=\min_{x\in\mathcal C}f(x)$.
\end{enumerate}

\section*{Main Convergence Theorem}
\begin{theorem}[Primal gap $O(1/k)$ for the classic step size]\label{thm:fw}
Assume \ref{A1}–\ref{A4} and use the step size
\[
\gamma_{k}=\frac{2}{k+2},\qquad k\ge0 .
\]
Then for all $k\ge0$
\[
f(x_{k})-f(x^{*})\;\le\;\frac{2L D^{2}}{k+2}.
\tag{3}
\]
\end{theorem}

\section*{Theoretical Properties}
\begin{enumerate}[label=\textbf{P\arabic*.},leftmargin=*]
\item \textbf{Stability.} Because each iterate is a convex combination of points in $\mathcal C$, all $x_{k}$ remain inside $\mathcal C$ (feasibility is preserved automatically).

\item \textbf{Hyper‑parameter sensitivity.}  
The classic schedule $\gamma_{k}=2/(k+2)$ is parameter‑free and yields the optimal $O(1/k)$ rate under the above assumptions. Fixed $\gamma$ may lead to a linear convergence plateau once the iterates approach the optimum.

\item \textbf{Comparison to projected gradient descent.}  
Projected GD needs a Euclidean projection onto $\mathcal C$ at each iteration, while Frank–Wolfe only requires solving the linear subproblem \eqref{1}. When the LMO is cheap (e.g., $\mathcal C$ is a simplex or a trace‑norm ball), Frank–Wolfe is computationally advantageous despite the slower $O(1/k)$ rate.

\item \textbf{Special case: Polytope with bounded curvature.}  
If the curvature constant $C_{f}$ (defined below) satisfies $C_{f}\le LD^{2}$, bound \eqref{3} can be tightened to $f(x_{k})-f(x^{*})\le \frac{2C_{f}}{k+2}$.
\end{enumerate}

\section*{Preliminaries (Definitions and Lemmas)}
\begin{definition}[Curvature constant]\label{def:curvature}
For a convex and $L$‑smooth $f$ on $\mathcal C$, define
\[
C_{f}\;:=\;\sup_{\substack{x,s\in\mathcal C\\ \gamma\in[0,1]}}
\frac{2}{\gamma^{2}}\Bigl(f\big(x+\gamma(s-x)\big)-f(x)-\gamma\langle\nabla f(x),s-x\rangle\Bigr).
\tag{4}
\]
\end{definition}
\begin{lemma}\label{lem:smoothness}
If $f$ is $L$‑smooth, then for any $x,s\in\mathcal C$ and $\gamma\in[0,1]$
\[
f\big(x+\gamma(s-x)\big)\le f(x)+\gamma\langle\nabla f(x),s-x\rangle
+\frac{L\gamma^{2}}{2}\|s-x\|^{2}.
\tag{5}
\]
\end{lemma}
\begin{proof}
From $L$‑smoothness we have  
$f(y)\le f(x)+\langle\nabla f(x),y-x\rangle+\frac{L}{2}\|y-x\|^{2}$ for any $x,y$.  
Set $y=x+\gamma(s-x)$ and obtain \eqref{5}. ∎
\end{proof}

\section*{Main Proof (Step‑by‑Step Derivation)}
We prove Theorem~\ref{thm:fw}. Throughout the proof let $x^{*}\in\arg\min_{\mathcal C}f$.

\begin{enumerate}[label={\bf [Step \arabic*]}]
\item \textbf{Descent inequality from the LMO.}  
By definition of $s_{k}$,
\[
\langle\nabla f(x_{k}),s_{k}-x_{k}\rangle
\;=\;\min_{s\in\mathcal C}\langle\nabla f(x_{k}),s-x_{k}\rangle
\;\le\;\langle\nabla f(x_{k}),x^{*}-x_{k}\rangle .
\tag{6}
\]

\item \textbf{Apply smoothness Lemma~\ref{lem:smoothness}.}  
With $y=x_{k+1}=x_{k}+\gamma_{k}(s_{k}-x_{k})$ and using \eqref{5},
\[
f(x_{k+1})\le f(x_{k})+\gamma_{k}\langle\nabla f(x_{k}),s_{k}-x_{k}\rangle
+\frac{L\gamma_{k}^{2}}{2}\|s_{k}-x_{k}\|^{2}.
\tag{7}
\]

\item \textbf{Bound the distance term.}  
Since $s_{k},x_{k}\in\mathcal C$ and $\operatorname{diam}(\mathcal C)=D$,  
\[
\|s_{k}-x_{k}\|\le D\quad\Longrightarrow\quad 
\|s_{k}-x_{k}\|^{2}\le D^{2}.
\tag{8}
\]

\item \textbf{Combine (6)–(8).}  
Insert the inequality from Step~1 into \eqref{7} and use \eqref{8}:
\[
f(x_{k+1})\le f(x_{k})+\gamma_{k}\langle\nabla f(x_{k}),x^{*}-x_{k}\rangle
+\frac{L D^{2}}{2}\gamma_{k}^{2}.
\tag{9}
\]

\item \textbf{Use convexity of $f$.}  
Convexity (\ref{A1}) yields
\[
f(x^{*})\ge f(x_{k})+\langle\nabla f(x_{k}),x^{*}-x_{k}\rangle .
\tag{10}
\]
Rearranging \eqref{10} gives
\[
\langle\nabla f(x_{k}),x^{*}-x_{k}\rangle\le f(x^{*})-f(x_{k}).
\tag{11}
\]

\item \textbf{Plug (11) into (9).}  
\[
f(x_{k+1})-f(x^{*})\le (1-\gamma_{k})\bigl[f(x_{k})-f(x^{*})\bigr]
+\frac{L D^{2}}{2}\gamma_{k}^{2}.
\tag{12}
\]

\item \textbf{Define the primal gap.}  
Let $h_{k}:=f(x_{k})-f(x^{*})\ge0$.  Equation \eqref{12} becomes
\[
h_{k+1}\le (1-\gamma_{k})h_{k}+\frac{L D^{2}}{2}\gamma_{k}^{2}.
\tag{13}
\]

\item \textbf{Insert the classic step size.}  
Set $\gamma_{k}=2/(k+2)$. Then $1-\gamma_{k}=(k)/(k+2)$. Substituting into \eqref{13}:
\[
h_{k+1}\le \frac{k}{k+2}\,h_{k}+ \frac{L D^{2}}{2}\Bigl(\frac{2}{k+2}\Bigr)^{2}
= \frac{k}{k+2}\,h_{k}+ \frac{2L D^{2}}{(k+2)^{2}}.
\tag{14}
\]

\item \textbf{Induction to obtain a closed form.}  
We prove by induction that for all $k\ge0$
\[
h_{k}\le \frac{2L D^{2}}{k+2}.
\tag{15}
\]

\begin{description}
\item[Base case $k=0$:]  
From the definition $h_{0}=f(x_{0})-f(x^{*})\le L D^{2}$ (by smoothness and compactness), thus
$h_{0}\le 2L D^{2}/2$, so \eqref{15} holds.

\item[Inductive step:]  
Assume \eqref{15} holds for $k$. Using \eqref{14}:
\[
h_{k+1}\le \frac{k}{k+2}\cdot\frac{2L D^{2}}{k+2}
+\frac{2L D^{2}}{(k+2)^{2}}
= \frac{2L D^{2}}{(k+2)^{2}}\bigl(k+1\bigr)
= \frac{2L D^{2}}{k+3}.
\]
Since $k+3=(k+1)+2$, the right‑hand side equals $2L D^{2}/(k+1+2)$, i.e. the bound \eqref{15} with $k\leftarrow k+1$. ∎
\end{description}

\item \textbf{Conclusion.}  
Inequality \eqref{15} is precisely the statement \eqref{3} of Theorem~\ref{thm:fw}. ∎
\end{enumerate}

\section*{Rate Derivation (With Constants)}
The bound \eqref{3} can be expressed with the curvature constant $C_{f}$ (Definition~\ref{def:curvature}).  
From Lemma~\ref{lem:smoothness} and the definition of $C_{f}$ we have $C_{f}\le L D^{2}$.  
Repeating the proof with $C_{f}$ in place of $L D^{2}$ yields
\[
f(x_{k})-f(x^{*})\le \frac{2 C_{f}}{k+2}.
\tag{16}
\]
Thus the constant $2L D^{2}$ in \eqref{3} is a (possibly loose) upper bound on $2C_{f}$.

\section*{Special Cases}
\begin{enumerate}[label=\textbf{S\arabic*.}]
\item \textbf{Polytope with vertices as extreme points.}  
When $\mathcal C$ is a polytope, each $s_{k}$ is an extreme vertex. The algorithm then produces a sparse representation of $x_{k}$ as a convex combination of at most $k+1$ vertices.

\item \textbf{Strongly convex $f$.}  
If $f$ is $\mu$‑strongly convex, Frank–Wolfe with exact line‑search enjoys a linear convergence rate:
\[
f(x_{k})-f(x^{*})\le \bigl(1-\tfrac{\mu}{L}\bigr)^{k}\bigl[f(x_{0})-f(x^{*})\bigr],
\]
provided the feasible set satisfies a suitable \emph{condition number} (e.g., a bounded pyramidal width).

\item \textbf{Non‑Euclidean norms.}  
The analysis extends verbatim to any norm $\|\cdot\|$ with corresponding smoothness constant $L_{\|\cdot\|}$ and diameter $D_{\|\cdot\|}$, yielding the same $O(1/k)$ rate.
\end{enumerate}

\newpage
\section*{Discussion}
The Frank–Wolfe method trades off a modest $O(1/k)$ convergence rate for the advantage of avoiding potentially expensive Euclidean projections. The proof above shows that the rate stems directly from smoothness, convexity, and the geometry of $\mathcal C$ (captured by $D$ or $C_{f}$). The step‑size $\gamma_{k}=2/(k+2)$ is the only hyper‑parameter required to achieve the optimal sublinear rate under the stated assumptions. Variants such as away‑steps, pairwise steps, or line‑search can improve the constant or even achieve linear convergence under stronger structural conditions (e.g., polyhedral sets and strong convexity).

\end{document}